% ============================================================================
% sec3_structural_model.tex — Compact framework (full model in Appendix A)
% ============================================================================

\section{Conceptual Framework}
\label{sec:structural}

A good cover bidder is invisible to the tools designed to catch
him. If he coordinates his bid tightly enough with the cartel's
designated winner, variance-based screens read the tender as
competitive. This section formalizes that intuition: when does
coordination defeat detection, and what trace does the cover bidder
leave behind? The answer generates five testable implications. Full
assumptions, proofs, and the structural likelihood are in
Appendices~\ref{app:proofs} and~\ref{app:structural_id}.

Consider a cartel that controls a designated winner and deploys
$m \geq 0$ cover bidders (FL firms), each bidding above the winning
bid $b^*$. The optimal $m^*$ falls with detection probability
$\theta_k$ and per-bidder cost $c_1$. Under convite, the
minimum-bidder rule ($\underline{n} = 3$) can force participation
even when market conditions would not justify it, mixing mandatory
and voluntary deployments and diluting the FL--price signal
(Section~\ref{sec:mechanisms}).

Two information structures generate two cover-bidding regimes.
Under \textbf{Regime~1} (complementary), cover bids scatter uniformly
above $b^*$, raising within-tender dispersion---and handing
variance-based screens an easy target. Under \textbf{Regime~2}
(coordinated), cover bids follow a truncated normal centered near
$b^*$, \emph{lowering} dispersion. This is the paradox at the heart
of the paper: the better the cartel coordinates, the more its
tenders \emph{look competitive} to dispersion-based tools. A
participation-based screen sidesteps the trap entirely---it does not
care how cover bids are distributed, only that cover bidders must
show up.

The model permits \textbf{strategic complementarity}: if the
marginal return to cover bidding is higher in competitive tenders
(many genuine bidders), then $m^*$ \emph{increases} in $n$. The
calibrated model confirms this pattern ($\hat{\gamma} = 0.69 > 0$;
Figure~\ref{fig:corner_solution} and
Appendix~\ref{app:structural_id}), indicating that cartels deploy
more cover bidders precisely where genuine competition is strongest.
This generates a \textbf{market-selection implication}: FL presence
should concentrate in competitive markets (low HHI).

Table~\ref{tab:predictions} maps the framework's diagnostic
implications to the empirical tests reported below.

\begin{table}[htbp]
\centering
\caption{Diagnostic implications and empirical tests}
\label{tab:predictions}
\small
\begin{tabular}{llll}
\toprule
\# & Diagnostic implication & Diagnostic test & Section \\
\midrule
D1\label{pred:price_v6} & FL associated with higher prices & Conditional price comparison & \ref{sec:reduced_form} \\
D2\label{pred:selection_v6} & FL adds to, not displaces, genuine bidders & Non-FL bidder count & \ref{sec:reduced_form} \\
D3\label{pred:regime_v6} & Coordinated regime: $\sigma_c < \sigma_g$ & BIC model selection & \ref{sec:structural_estimation} \\
D4\label{pred:market_v6} & FL association $\downarrow$ in HHI & Network split & \ref{sec:empirical_network} \\
D5\label{pred:exchange_v6} & FL residuals non-exchangeable & Bajari--Ye KS + pairwise & \ref{sec:bajari_ye} \\
\bottomrule
\end{tabular}
\end{table}

\begin{figure}[htbp]
\centering
\includegraphics[width=0.95\textwidth]{fig_corner_solution.pdf}
\caption{Optimal number of cover bidders $m^*$ as a function of genuine
bidders $n$. Dotted region ($n < 3$): constraint-binding corner
solution under the minimum-bidder rule ($\underline{n} = 3$,
Lei~8.666/93). Solid curve: calibrated interior solution
($\hat{\gamma} = 0.69$, strategic complementarity). Points: empirical
binned means with 95\% CIs.}
\label{fig:corner_solution}
\end{figure}

\noindent These are predictions, not conclusions. The next sections
confront them with data: does FL presence actually predict higher
prices? Do FL bids really cluster where the framework says they
should? And does the screen find cartels that existing tools miss?
Illustrative welfare calculations appear in
Appendix~\ref{app:extensions}.
